% Part: many-valued-logic
% Chapter: supervaluationist-logic
% Section: semantics

\documentclass[../../../include/open-logic-section]{subfiles}

\begin{document}

\olfileid{mvl}{sup}{sem}

\olsection{Semantics}

Supervaluationist has the same !!{structure}s as three-valued logics and (negative) free logics. 

\begin{defn}[!!^{structure}s] 
A \emph{!!{valuation}} for a propositional language $\Lang{L}$ is a function~$\pAssign{v}$ assigning $\True$, $\False$ or $\Undef$ to the !!{propositional variable}s of the language, i.e., $\pAssign{v} \colon
\PVar \to {\True,\False,\Undef}$.

A \emph{!!{structure}} for a first-order language $\Lang{L}$ is a triple $\Struct{M}$ consisting of !!a{domain} of objects, $\Domain{M}$, and an interpretation $\Assign{\cdots}{M}$ of constants and predicates that can be partially defined.
\end{defn}

(When applying supervaluationism to vagueness we also allow predicates to be partially defined for objects in the domain. That is, instead of the interpretation of a predicate being a set of tuples, we make it a function that maps tuples to $\True$, $\False$ or $\Undef$.)

The difference lies in the semantics for complex !!{sentence}s. Given !!a{structure}, we consider all the ways in which we could turn its non-classical elements into classical values. We call these \emph{supervaluations}.

\begin{defn}[Supervaluations]
A \emph{supervaluation} $\pAssign{v^+}$ for $\pAssign{v}$ is a function that maps every !!{propositional variable} to $\True$ or $\False$, such that $\pAssign{v^+}(!A)=\pAssign{v}(!A)$ if $\pAssign{v}(!A) = \True$ or $\False$. 

A \emph{supervaluation} $\Struct{M^+}$ of $\Struct{M}$ is !!a{structure} such that $\Domain{M^+}=\Domain{M}$ and $\Assign{\cdot}{M^+}$ is a classical interpretation such that $\Assign{c}{M^+} = \Assign{c}{M^+}$ if the latter is defined, and $\Assign{R}{M^+} = \Assign{R}{M}$ if the latter is defined.
\end{defn}

(When applying vagueness to supervaluationism, we say where $R$ is a $n$-place predicate, $\Assign{R}{M^+}$ assigns $\True$ or $\False$ to every tuple and agrees with $\Assign{R}{M}$ whenever the latter assigns $\True$ or $\False$ to some tuple.)

Note that supervaluations are \emph{classical} !!{structure}s. If you think of $\Undef$, empty names and predicates with partial extension as being (partially or wholly) undefined, supervaluations are ways of making them defined (classical). Supervaluationist semantics, then, is simply the idea that !!a{sentence} has a classical truth-value only if \emph{all} the ways of making it defined (classical) have that truth-value, and is neither truth nor false otherwise.

\begin{defn}[Supervaluationist semantics]
Truth and satisfaction of a !!{formula}~$!A$ \emph{at a supervaluation}, $\pAssign{v^+}(A)$, $\Sat{M^+}{!A}$ is defined as in classical logic.

The value of !!a{sentence} $!A$ at a !!a{valuation} $\pAssign{v}$ is $\True$ if $\pAssign{v^+}(!A) = \True$ for every supervaluation $\pAssign{v^+}$ for $\pAssign{v}$, $\False$ if $\pAssign{v^+}(!A) = \False$ for every supervaluation $\pAssign{v^+}$ for $\pAssign{v}$, and $\Undef$ otherwise.

Analogously, the value of !!a{sentence} $!A$ at !!a{structure} $\Struct{M}$,  is $\Value{!A}{M}$, is $\True$ if $\Value{M^+}(!A) = \True$ for every supervaluation $\Struct{M^+}$ for $\Struct{M}$, $\False$ if $\pAssign{v^+}(!A) = \False$ for every supervaluation $\Struct{M^+}$ for $\Struct{M}$, and $\Undef$ otherwise.
\end{defn}

!!^{sentence}s that are true are every supervaluation or false at every supervaluation are sometimes often called `supertrue' and `superfalse'. In the traditional supervaluationist picture, these are equated with truth and falsehood. 

Using $\True$ as the designated truth value, the logic of supervaluationism turns out to be classical logic.

\begin{prop}
$\Gamma \Entails[\LogSup] !A$ iff $\Gamma \Entails[\LogCL] !A$
\end{prop}

This holds only as long as the language is \emph{not} supplemented with operators to talk about (super)truth, as we see below.

Some obvious but noteworthy consquences:

\begin{prop}
\begin{itemize}
	\item $\Entails[\LogSup] !A \lor \lnot !A$
	\item $\Entails[\LogSup] \lnot (!A \lor \lnot !A)$
	\item For any sub-classical logic $\Log{L}$ i(e.g. $\LogLuk, \LogKw, \LogKs, \ldots$), if $\Gamma \Entails[\Log{L}] !A$ then $\Gamma \Entails[\LogSup] !A$.
\end{itemize}
\end{prop}
	
\end{document}
